\documentclass[12pt]{article}
\usepackage[utf8]{inputenc}
\usepackage{amsmath}
\usepackage{amssymb}
\usepackage{graphicx}
\usepackage{hyperref}
\usepackage{geometry}
\geometry{a4paper, margin=1in}
\usepackage{xcolor}
\usepackage{listings}

\title{\textbf{Deconstructing RWKV \\ Part 2: Pure PyTorch Implementation}}
\author{Sovesh Mohapatra}
\date{\today}

\begin{document}

\maketitle

\section*{Introduction}
In Part 1 of this series, we explored the mathematical foundations of RWKV: how attention can be rewritten as a linear recurrence via the WKV operator, enabling parallel training and $O(1)$ inference.

Today, in Part 2, we turn that math into code.

I've written a complete RWKV implementation in pure PyTorch—no external libraries, no CUDA kernels, just clean tensor operations. The code supports both parallel training mode (for efficiency) and recurrent inference mode (for memory-flat generation).

\section*{The Implementation Strategy}

RWKV's elegance lies in its simplicity. The architecture decomposes into two modules:

\begin{enumerate}
    \item \textbf{Time Mixing (TM)}: The attention equivalent, computing the WKV recurrence.
    \item \textbf{Channel Mixing (CM)}: The FFN equivalent, using gated squared ReLU.
\end{enumerate}

Each module follows a pre-norm pattern similar to Transformers:
\begin{equation}
    x \leftarrow x + \text{Module}(\text{LayerNorm}(x))
\end{equation}

\section*{Time Mixing: Parallel Training Mode}

The key challenge in implementing RWKV is computing the WKV operator efficiently. During training, we have access to the entire sequence at once, so we want to parallelize the recurrence.

Recall the WKV recurrence from Part 1:
\begin{align}
    \text{num}_t &= \exp(u + k_t) v_t + \exp(-w) \cdot \text{num}_{t-1} \\
    \text{den}_t &= \exp(u + k_t) + \exp(-w) \cdot \text{den}_{t-1} \\
    \text{wkv}_t &= \frac{\text{num}_t}{\text{den}_t}
\end{align}

For parallel computation, we can unroll this recurrence. The numerator at time $t$ is:
\begin{equation}
    \text{num}_t = \exp(u + k_t) v_t + \sum_{i=0}^{t-1} \exp(-(t-1-i)w + k_i) v_i
\end{equation}

This can be computed using a \textbf{cumulative sum with decay}:
\begin{enumerate}
    \item Compute $kv_t = k_t \cdot v_t$ for all $t$.
    \item Compute decay powers: $\text{decay}^t = \exp(-tw)$.
    \item Weight each $kv_t$ by its decay: $\text{weighted}_t = kv_t \cdot \text{decay}^t$.
    \item Cumulative sum: $\text{cumsum}_t = \sum_{i=0}^t \text{weighted}_i$.
    \item Rescale to get proper numerator.
\end{enumerate}

In PyTorch, this becomes:
\begin{lstlisting}[language=Python, basicstyle=\small\ttfamily]
# Compute weighted kv
kv = k * v  # element-wise
decay_powers = torch.pow(decay, time_indices)
weighted_kv = kv * decay_powers

# Cumulative sum gives the recurrence
cumsum = torch.cumsum(weighted_kv, dim=1)

# First token uses time_first, rest use decay
wkv[:, 0] = time_first * kv[:, 0]
wkv[:, 1:] = time_first * kv[:, 1:] + decay * cumsum[:, :-1]
\end{lstlisting}

\section*{Time Mixing: Recurrent Inference Mode}

During inference, we generate tokens one at a time. Here, the recurrence shines—we simply maintain the accumulated state:

\begin{lstlisting}[language=Python, basicstyle=\small\ttfamily]
def forward_recurrent(self, x_t, state):
    # Project
    k_t = self.key(x_t)
    v_t = self.value(x_t)
    r_t = self.receptance(x_t)
    
    # WKV recurrence (single step)
    wkv_t = time_first * k_t * v_t + decay * state
    
    # Apply receptance gate
    output = sigmoid(r_t) * wkv_t
    
    # Return output and new state
    return output, wkv_t
\end{lstlisting}

The state is just a single tensor of shape $(B, C)$—constant size regardless of how many tokens we've generated. Compare this to Transformers, which must store $(B, L, C)$ for the KV cache!

\section*{Channel Mixing: Squared ReLU Gating}

The Channel Mixing module is simpler. It uses a GLU-style gating mechanism with squared ReLU:

\begin{lstlisting}[language=Python, basicstyle=\small\ttfamily]
def forward(self, x):
    x = self.ln(x)  # pre-norm
    
    # Key projection with squared ReLU
    k = self.key(x)
    k_squared = F.relu(k) ** 2
    
    # Receptance gate
    r = torch.sigmoid(self.receptance(x))
    
    # Value projection and gate
    v = self.value(k_squared)
    output = r * v
    
    return output
\end{lstlisting}

The squared ReLU ($\text{ReLU}(x)^2$) provides smoother gradients than standard ReLU while maintaining sparsity. This is a subtle but important design choice that improves training stability.

\section*{The Complete RWKV Model}

Stacking Time Mixing and Channel Mixing blocks gives us the full RWKV model:

\begin{lstlisting}[language=Python, basicstyle=\small\ttfamily]
class RWKV(nn.Module):
    def __init__(self, vocab_size, embed_dim, num_layers):
        super().__init__()
        self.embedding = nn.Embedding(vocab_size, embed_dim)
        self.blocks = nn.ModuleList([
            RWKVBlock(embed_dim) for _ in range(num_layers)
        ])
        self.ln_out = nn.LayerNorm(embed_dim)
        self.head = nn.Linear(embed_dim, vocab_size, bias=False)
        
        # Weight tying (like Transformers)
        self.head.weight = self.embedding.weight
    
    def forward(self, x, state=None, use_recurrent=False):
        x = self.embedding(x)
        
        for block in self.blocks:
            x, state = block(x, state, use_recurrent)
        
        x = self.ln_out(x)
        return self.head(x), state
\end{lstlisting}

Weight tying between the embedding and output projection is a standard technique that improves parameter efficiency and training stability.

\section*{Autoregressive Generation}

With the recurrent mode implemented, autoregressive generation is straightforward:

\begin{lstlisting}[language=Python, basicstyle=\small\ttfamily]
@torch.no_grad()
def generate(self, prompt, max_new_tokens, temperature=1.0):
    self.eval()
    context = prompt.clone()
    state = None
    
    # Process prompt in parallel
    logits, state = self(context, use_recurrent=True)
    
    for _ in range(max_new_tokens):
        # Sample next token
        probs = F.softmax(logits / temperature, dim=-1)
        next_token = torch.multinomial(probs, 1)
        
        # Append and step
        context = torch.cat([context, next_token], dim=1)
        logits, state = self(next_token, state, use_recurrent=True)
    
    return context
\end{lstlisting}

The key insight: after processing the prompt in parallel, each new token requires only a single recurrent step with $O(1)$ memory.

\section*{Implementation Gotchas}

A few pitfalls I encountered while implementing RWKV:

\begin{enumerate}
    \item \textbf{Decay Parameterization}: The decay $w$ must be positive to ensure stable fading memory. I parameterize it via sigmoid to keep it in $(0, 1)$.
    
    \item \textbf{Time First Initialization}: The $u$ parameter (time\_first) should be initialized to emphasize the current token. I use $u = 1.0$ (giving $\sigma(u) \approx 0.73$).
    
    \item \textbf{Gradient Clipping}: RWKV can have gradient explosions in early training. Gradient clipping at norm 1.0 is essential.
    
    \item \textbf{State Initialization}: During recurrent inference, the state starts at zero. This represents no prior context.
\end{enumerate}

\section*{Next Steps: Benchmarking}

The code is complete. In Part 3, we'll put RWKV to the test against Transformer and LSTM baselines on synthetic sequence tasks, measuring training convergence, inference latency, and memory usage.

\vspace{1em}
\noindent \textit{The full PyTorch implementation is available on GitHub. Stay tuned for the benchmark results!}

\end{document}
